\documentclass[12pt,a4paper]{article}

\usepackage{fontawesome5}
\providecommand{\faPencilAlt}{\faIcon{pencil-alt}}

\usepackage{styles/couleurs}
\usepackage{styles/mise_en_page}
\usepackage{styles/boites}
\usepackage{styles/maths}
\usepackage{styles/figures}

\usepackage{enumitem}
\usepackage{array}
\usepackage{tabularx}
\usepackage{amsmath}
\usepackage{amssymb}
\usepackage{graphicx}
\usepackage{tikz}
\usepackage{xparse}

\usetikzlibrary{calc,angles,quotes,arrows.meta}

\setlist[itemize]{leftmargin=1.4em,itemsep=0.15em,topsep=0.2em}
\setlist[enumerate]{leftmargin=*,itemsep=0.22em,topsep=0.25em}

\newcommand{\lignepointillee}{\noindent\dotfill\par\vspace{0.45em}}

\newcommand{\lignesreponse}[1]{%
\begin{tcolorbox}[
enhanced,colback=white,colframe=NeutreC!40,boxrule=0.5pt,
arc=1.5mm,left=3mm,right=3mm,top=2mm,bottom=1mm,
before skip=2pt,after skip=2pt]
\foreach \i in {1,...,#1}{\lignepointillee}
\end{tcolorbox}}

\newcommand{\titreseq}[1]{%
\begin{center}
{\sffamily\LARGE\bfseries\color{TitreC}Séquence 6 -- Cosinus d'un angle aigu}\\[0.2em]
{\sffamily\large\bfseries\color{SavoirC}#1}
\end{center}}

\begin{document}

\titreseq{Cours -- Calculer des longueurs dans un triangle rectangle}

\section{Vocabulaire du triangle rectangle}

\definition{Dans un triangle rectangle :

\begin{itemize}
  \item l'\textbf{hypoténuse} est le côté opposé à l'angle droit ;
  \item c'est aussi le plus long côté du triangle ;
  \item pour un angle aigu donné, le \textbf{côté adjacent} est le côté de l'angle droit
        qui touche cet angle.
\end{itemize}}

\begin{center}
\begin{tikzpicture}[scale=0.9]
  \coordinate (A) at (0,0);
  \coordinate (B) at (5.5,0);
  \coordinate (C) at (5.5,3.8);
  \draw[line width=0.9pt] (A)--(B)--(C)--cycle;
  \draw (5.15,0)--(5.15,0.35)--(5.5,0.35);
  \draw pic[draw=ActionC,fill=ActionC!15,angle radius=8mm] {angle=B--A--C};
  \node[below left] at (A) {$A$};
  \node[below right] at (B) {$B$};
  \node[above right] at (C) {$C$};
  \node[below] at ($(A)!0.5!(B)$) {côté adjacent à \(\widehat A\)};
  \node[above left] at ($(A)!0.55!(C)$) {hypoténuse};
\end{tikzpicture}
\end{center}

\remarque{Le côté adjacent dépend de l'angle étudié.
Dans le même triangle, le côté adjacent à \(\widehat A\) n'est pas
le côté adjacent à \(\widehat C\).}

\section{Définition du cosinus}

\definition{Dans un triangle rectangle, le cosinus d'un angle aigu est le quotient :

\[
\cos(\text{angle})
=
\text{côté adjacent}\div\text{hypoténuse}.
\]}

\exemple{Dans le triangle \(ABC\) rectangle en \(B\) :

\[
\cos(\widehat A)=AB\div AC
\]

et

\[
\cos(\widehat C)=BC\div AC.
\]}

\propriete{Le cosinus d'un angle aigu est un nombre sans unité compris entre \(0\) et \(1\).}

\remarque{Le cosinus dépend de l'angle, mais pas de la taille du triangle.
Des triangles rectangles ayant le même angle aigu donnent le même quotient
« côté adjacent divisé par hypoténuse ».}

\begin{center}
\begin{minipage}[t]{0.31\textwidth}
\centering
\textbf{Famille 1 : \(35^\circ\)}\par\medskip
\begin{tikzpicture}[scale=0.62]
  \coordinate (O) at (0,0);
  \draw[line width=0.9pt] (O)--(7.7,0);
  \draw[line width=0.9pt] (O)--({8*cos(35)},{8*sin(35)});
  \foreach \r/\name in {4/A,6/B,8/C}{
    \coordinate (\name) at ({\r*cos(35)},{\r*sin(35)});
    \coordinate (\name-p) at ({\r*cos(35)},0);
    \draw (\name)--(\name-p);
    \draw ($(\name-p)+(-0.28,0)$)--($(\name-p)+(-0.28,0.28)$)--($(\name-p)+(0,0.28)$);
  }
  \draw pic[draw=ActionC,angle radius=8mm,"$35^\circ$"] {angle=C-p--O--C};
\end{tikzpicture}
\end{minipage}
\hfill
\begin{minipage}[t]{0.31\textwidth}
\centering
\textbf{Famille 2 : \(50^\circ\)}\par\medskip
\begin{tikzpicture}[scale=0.62]
  \coordinate (O) at (0,0);
  \draw[line width=0.9pt] (O)--(5.3,0);
  \draw[line width=0.9pt] (O)--({8*cos(50)},{8*sin(50)});
  \foreach \r/\name in {4/D,6/E,8/F}{
    \coordinate (\name) at ({\r*cos(50)},{\r*sin(50)});
    \coordinate (\name-p) at ({\r*cos(50)},0);
    \draw (\name)--(\name-p);
    \draw ($(\name-p)+(-0.28,0)$)--($(\name-p)+(-0.28,0.28)$)--($(\name-p)+(0,0.28)$);
  }
  \draw pic[draw=ActionC,angle radius=8mm,"$50^\circ$"] {angle=F-p--O--F};
\end{tikzpicture}
\end{minipage}
\hfill
\begin{minipage}[t]{0.31\textwidth}
\centering
\textbf{Famille 3 : \(70^\circ\)}\par\medskip
\begin{tikzpicture}[scale=0.62]
  \coordinate (O) at (0,0);
  \draw[line width=0.9pt] (O)--(3,0);
  \draw[line width=0.9pt] (O)--({8*cos(70)},{8*sin(70)});
  \foreach \r/\name in {4/G,6/H,8/I}{
    \coordinate (\name) at ({\r*cos(70)},{\r*sin(70)});
    \coordinate (\name-p) at ({\r*cos(70)},0);
    \draw (\name)--(\name-p);
    \draw ($(\name-p)+(-0.28,0)$)--($(\name-p)+(-0.28,0.28)$)--($(\name-p)+(0,0.28)$);
  }
  \draw pic[draw=ActionC,angle radius=8mm,"$70^\circ$"] {angle=I-p--O--I};
\end{tikzpicture}
\end{minipage}
\end{center}


\section{Utiliser la calculatrice}

\methode{Avant tout calcul :

\begin{enumerate}[label=\textbf{\arabic*.}]
  \item vérifier que la calculatrice est en mode degré ;
  \item utiliser la touche \(\cos\) avec la mesure de l'angle ;
  \item conserver la valeur affichée jusqu'à la dernière étape ;
  \item arrondir seulement le résultat final.
\end{enumerate}}

\exemple{Contrôle rapide du mode degré :

\[
\cos 60^\circ=0{,}5.
\]

Si la calculatrice n'affiche pas \(0{,}5\), son réglage doit être vérifié.}

\newpage

\section{Calculer le côté adjacent}

\propriete{À partir de la définition :

\[
\cos(\text{angle})
=
\text{adjacent}\div\text{hypoténuse},
\]

on obtient :

\[
\boxed{
\text{côté adjacent}
=
\text{hypoténuse}\times\cos(\text{angle})
}
\]}

\methode{Pour calculer le côté adjacent :

\begin{enumerate}[label=\textbf{\arabic*.}]
  \item nommer le triangle rectangle ;
  \item repérer l'hypoténuse et le côté adjacent ;
  \item écrire la relation du cosinus avec les noms des segments ;
  \item multiplier l'hypoténuse par le cosinus de l'angle ;
  \item arrondir et écrire l'unité.
\end{enumerate}}

\exemple{Dans le triangle \(DEF\) rectangle en \(E\),
on connaît \(DF=12\) cm et \(\widehat D=38^\circ\).
On cherche \(DE\).

\[
\cos(\widehat D)=DE\div DF.
\]

Donc :

\[
DE=DF\times\cos(\widehat D).
\]

\[
DE=12\times\cos 38^\circ\approx9{,}5.
\]

Ainsi :

\[
\boxed{DE\approx9{,}5\text{ cm}}.
\]}

\newpage

\section{Calculer l'hypoténuse}

\propriete{Comme :

\[
\text{adjacent}
=
\text{hypoténuse}\times\cos(\text{angle}),
\]

on défait la multiplication en divisant :

\[
\boxed{
\text{hypoténuse}
=
\text{côté adjacent}\div\cos(\text{angle})
}
\]}

\methode{Pour calculer l'hypoténuse :

\begin{enumerate}[label=\textbf{\arabic*.}]
  \item nommer le triangle rectangle ;
  \item repérer l'hypoténuse et le côté adjacent ;
  \item écrire la relation du cosinus avec les noms des segments ;
  \item diviser le côté adjacent par le cosinus de l'angle ;
  \item arrondir et écrire l'unité.
\end{enumerate}}

\exemple{Dans le triangle \(GHI\) rectangle en \(H\),
on connaît \(GH=7\) cm et \(\widehat G=41^\circ\).
On cherche \(GI\).

\[
\cos(\widehat G)=GH\div GI.
\]

Donc :

\[
GI=GH\div\cos(\widehat G).
\]

\[
GI=7\div\cos 41^\circ\approx9{,}3.
\]

Ainsi :

\[
\boxed{GI\approx9{,}3\text{ cm}}.
\]}

\newpage

\section{Choisir entre multiplier et diviser}

\begin{center}
\renewcommand{\arraystretch}{1.5}
\begin{tabularx}{0.9\textwidth}{
|>{\raggedright\arraybackslash}X
|>{\centering\arraybackslash}p{5cm}|}
\hline
\textbf{Longueur cherchée} & \textbf{Calcul}\\
\hline
Côté adjacent &
\(\text{hypoténuse}\times\cos(\text{angle})\)\\
\hline
Hypoténuse &
\(\text{côté adjacent}\div\cos(\text{angle})\)\\
\hline
\end{tabularx}
\end{center}

\remarque{Il est préférable de repartir de la phrase :

\[
\text{adjacent}
=
\text{hypoténuse}\times\cos(\text{angle})
\]

plutôt que de mémoriser deux formules sans les comprendre.}

\section{Rédiger et contrôler}

\methode{Une rédaction complète contient :

\begin{enumerate}[label=\textbf{\arabic*.}]
  \item le nom du triangle et le sommet de l'angle droit ;
  \item la relation du cosinus avec les noms des segments ;
  \item l'égalité donnant la longueur cherchée ;
  \item le calcul à la calculatrice ;
  \item l'arrondi, l'unité et une phrase de conclusion.
\end{enumerate}}

\propriete{Dans un triangle rectangle :

\begin{itemize}
  \item l'hypoténuse est plus longue que chacun des deux autres côtés ;
  \item le côté adjacent est donc plus petit que l'hypoténuse ;
  \item multiplier l'hypoténuse par un cosinus compris entre \(0\) et \(1\)
        donne une longueur plus petite ;
  \item diviser le côté adjacent par un cosinus compris entre \(0\) et \(1\)
        donne une longueur plus grande.
\end{itemize}}

\erreur{Erreurs fréquentes :

\begin{itemize}
  \item utiliser le cosinus dans un triangle qui n'est pas rectangle ;
  \item choisir un côté adjacent sans préciser l'angle étudié ;
  \item confondre côté adjacent et hypoténuse ;
  \item multiplier alors que l'on cherche l'hypoténuse ;
  \item utiliser la calculatrice en mode radian ;
  \item arrondir les valeurs intermédiaires ;
  \item oublier l'unité ou obtenir un côté adjacent plus long que l'hypoténuse.
\end{itemize}}

\section{Dans un problème}

\methode{Pour résoudre un problème :

\begin{enumerate}[label=\textbf{\arabic*.}]
  \item faire un schéma simple et codé ;
  \item repérer ou construire le triangle rectangle utile ;
  \item traduire les données dans ce triangle ;
  \item choisir le calcul de cosinus adapté ;
  \item utiliser, si nécessaire, une autre propriété ensuite,
        par exemple le théorème de Pythagore ;
  \item interpréter le résultat dans la situation.
\end{enumerate}}

%==================================================
% Encart calculatrice : mode degré
%==================================================

\begin{tcolorbox}[
  enhanced,
  colback=SavoirBg,
  colframe=SavoirC!65,
  boxrule=0.7pt,
  arc=2mm,
  left=3mm,
  right=3mm,
  top=2mm,
  bottom=2mm,
  title={\sffamily\bfseries
    \faCalculator\ Calculatrice : vérifier le mode degré},
  coltitle=SavoirC,
  fonttitle=\sffamily\bfseries
]

Avant d'effectuer un calcul avec un cosinus, vérifie que la calculatrice
est réglée en \textbf{degrés}.

\vspace{0.3em}

\renewcommand{\arraystretch}{1.35}
\begin{tabularx}{\linewidth}{
  @{}
  >{\sffamily\bfseries\color{TitreC}}p{3.2cm}
  X
  @{}
}
\textbf{Casio fx-92}
&
\texttt{CONFIG}
\(\longrightarrow\)
\texttt{Unité d'angle}
\(\longrightarrow\)
\texttt{Degré}
\\

\textbf{TI-Collège Plus}
&
\texttt{MODE}
\(\longrightarrow\)
sélectionner \texttt{DEG}
\(\longrightarrow\)
valider
\end{tabularx}

\begin{center}
  \sffamily\bfseries
  L'écran doit afficher
  \tcbox[
    on line,
    colback=white,
    colframe=SavoirC,
    boxrule=0.6pt,
    arc=1mm,
    left=1.5mm,
    right=1.5mm,
    top=0.4mm,
    bottom=0.4mm
  ]{\texttt{D} ou \texttt{DEG}}.
\end{center}

\textbf{Test rapide :}

\[
\cos(60^\circ)=0{,}5.
\]

Si la calculatrice n'affiche pas \(0{,}5\), vérifie à nouveau le mode.

\end{tcolorbox}


\aretenir{\textbf{La phrase essentielle :}

\[
\boxed{
\text{côté adjacent}
=
\text{hypoténuse}\times\cos(\text{angle aigu})
}
\]

Le cosinus ne s'utilise que dans un triangle rectangle.
Le côté adjacent dépend de l'angle étudié.
On contrôle toujours le mode degré, la cohérence du résultat, l'arrondi et l'unité.}

\end{document}
